\chapter{Introdução} \label{cap:introducao}

A Internet das Coisas (IoT) é composta por dispositivos embarcados que desempenham funções específicas, atuando como sensores e controladores para uma variedade de aplicações tecnológicas produzidas hoje em dia, desde os aparelhos mais simples do cotidiano até processos industriais complexos. Com a evolução do hardware, os dispositivos IoT tornaram-se menores, mais rápidos, mais eficientes e com menor custo, o que levou a uma explosão do seu uso \cite{alfuqaha2015iot}.

Contudo, um dos principais entraves para essa tecnologia decorre das limitações energéticas e computacionais desses dispositivos \cite{nistir8114}. Apesar da evolução vista nas últimas décadas, a IoT ainda enfrenta gargalos em processos que demandam alto poder computacional, devido à grande carga de processamento e ao elevado custo energético.

Diante desse cenário, surgem questões de segurança que podem ser exploradas para invadir e comprometer o funcionamento da rede. Para resolver a questão da transmissão segura e eficiente, padrões como o ASCON, pela necessidade de se padronizar e produzir um conjunto de métodos criptográficos que sejam leves para uso em dispositivos com baixo consumo energético \cite{nist800232}. Em paralelo, para viabilizar a comunicação de longa distância com baixa energia, utiliza-se amplamente o protocolo LoRaWAN, que possui aplicações fundamentais em diversas áreas, desde a infraestrutura de cidades inteligentes e controle de tráfego até o monitoramento de desastres, agricultura de precisão e automação industrial com foco em manutenção preditiva. O grande problema enfrentado por esse protocolo, no entanto, é a 
segurança da informação frente a ataques como jamming, replay attacks, 
spoofing e denial of service \cite{hessel2023lorawan}. A vulnerabilidade a ataques podem causar disrupções críticas e prejuízos severos, tornando indispensável a aplicação de métodos robustos de prevenção. 

%%%%%%%%%%%%%%%%%%%%%%%%%%%%%%%%%%%%%%%%%%%%%%%%%%%%
%%%%%%%%%%%%%%%%%%%%%%%%%%%%%%%%%%%%%%%%%%%%%%%%%%%%

\section{Motivação}
A necessidade de operar com baixo consumo energético eleva bastante a complexidade do uso de dispositivos IoT, uma vez que tentativas de mitigar esse problema podem resultar em simplificações que geram problemas de segurança \cite{nistir8114} \cite{hessel2023lorawan}. Diante disso, torna-se pertinente compreender a natureza desses ataques e suas formas de mitigação, além de estudar a eficiência e a eficácia de modelos de Machine Learning na detecção de intrusões e anomalias.

Um desses problemas é a vulnerabilidade a ataques diretos, como ataques de negação de serviço (DDoS) ou Jamming, que podem esgotar a bateria do dispositivo ou impedir que a coleta de dados seja feita de maneira adequada \cite{hessel2023lorawan}. Outro grande desafio está na comunicação: é preciso garantir que a informação transmitida entre os dispositivos não possa ser interceptada nem modificada, assegurando a integridade, confidencialidade e autenticidade dos dados. Como os algoritmos de segurança comumente utilizados exigem um grande poder computacional, torna-se necessário buscar maneiras mais eficientes de proteger os dados durante a transmissão.


O uso de \textit{Machine Learning} neste trabalho decorre da necessidade de 
detectar anomalias em redes LoRaWAN a partir de padrões temporais e 
comportamentais dos dados. Em contextos reais de IoT, abordagens baseadas 
apenas em regras fixas ou limiares estáticos tendem a ser limitadas. O 
comportamento normal pode variar conforme o contexto operacional (dispositivo, 
gateway, ambiente) e evoluir ao longo do tempo, tornando limiares fixos 
(ex: ``se RSSI < -120 dBm, alertar'') inadequados. Além disso, métodos 
tradicionais apresentam dificuldade em capturar estruturas complexas em dados 
sequenciais, tais como degradação gradual de equipamento ou padrões não 
lineares de transmissão. Por outro lado, abordagens baseadas em aprendizado 
conseguem modelar relações temporais, identificar desvios sutis e adaptar-se 
dinamicamente à evolução do comportamento normal, mostrando-se mais adequadas 
para este contexto \cite{chalapathy2019}.

Por isso, a modelagem é realizada com base no comportamento esperado do dispositivo em suas circunstâncias normais de operação, buscando aprender padrões legítimos de funcionamento a partir dos dados observados. Nessa abordagem, desvios suficientemente relevantes em relação ao comportamento previamente aprendido são tratados como indícios de anomalia. Trata-se, portanto, de uma abordagem adaptativa, pois o modelo se ajusta ao perfil comportamental de cada nó monitorado. \cite{hayashi2022handsfree}.


Assim, neste trabalho, \textit{Machine Learning} é tratado como um dos núcleos centrais da investigação. Seu papel não é apenas classificar dados, mas aprender perfis de operação, comparar abordagens de detecção e oferecer uma base experimental para uma etapa posterior de avaliação em borda. Essa decisão amplia o escopo da análise, permitindo estudar simultaneamente eficácia de detecção, robustez metodológica e viabilidade de implantação em borda.

O padrão Ascon, por sua vez, foi selecionado como a primeira escolha como novo padrão da competição de criptografia leve do NIST \citeonline{nist_lwc_finalists} assim como a primeira escolha de criptografia leve da competição CAESAR \cite{caesar_submissions}. A relevância desse padrão se torna clara quando analisado seu baixo custo energético, rapidez, e confiabilidade, garantindo segurança para dados AEAD e ataques de canal colateral \cite{ascon_official}.

%%%%%%%%%%%%%%%%%%%%%%%%%%%%%%%%%%%%%%%%%%%%%%%%%%%%
%%%%%%%%%%%%%%%%%%%%%%%%%%%%%%%%%%%%%%%%%%%%%%%%%%%%

\section{Objetivos}
O objetivo geral deste trabalho é implementar e avaliar um sistema de comunicação baseado no protocolo LoRaWAN, integrando a aplicação de modelos de \textit{Machine Learning} e o padrão de criptografia leve Ascon. A finalidade central é analisar o nível de segurança e a resiliência que a combinação dessas duas tecnologias pode oferecer ao ecossistema de Internet das Coisas (IoT).

Para o alcance do objetivo geral, foram definidos os seguintes objetivos específicos:

\begin{itemize}
    \item Treinar e testar modelos de \textit{Machine Learning}, utilizando os \textit{datasets} selecionados, para atuar na identificação e mitigação de ameaças ativas, tais como \textit{jamming}, \textit{spoofing} e ataques de Negação de Serviço (DoS).
    \item Implementar os algoritmos da família Ascon no nó sensor para estabelecer uma infraestrutura que garanta a privacidade, a autenticidade e a integridade dos dados transmitidos, com baixo custo computacional.
    \item Desenvolver o protótipo do sistema integrando as camadas de comunicação (LoRaWAN), de detecção e criptografia.
    \item Analisar o nível de segurança global oferecido pelo conjunto, validando a eficácia da solução proposta frente às restrições do hardware - em termos de consumo energético.
\end{itemize}

%%%%%%%%%%%%%%%%%%%%%%%%%%%%%%%%%%%%%%%%%%%%%%%%%%%%
%%%%%%%%%%%%%%%%%%%%%%%%%%%%%%%%%%%%%%%%%%%%%%%%%%%%

\section{Justificativa}
A relevância deste trabalho evidencia-se pelo papel central e pela rápida expansão dos dispositivos de Internet das Coisas (IoT) no cotidiano, com projeções indicando dezenas de bilhões de equipamentos conectados nos próximos anos \cite{pandey2025}. Propor uma metodologia que integra tecnologias e conceitos críticos como o protocolo de comunicação LoRaWAN, modelos de \textit{Machine Learning} e segurança criptográfica, e verificar a sua eficácia em cenários práticos, demonstra como a união dessas ferramentas pode alavancar a resiliência de sistemas reais, tornando-os cada vez mais robustos contra ameaças.

%%%%%%%%%%%%%%%%%%%%%%%%%%%%%%%%%%%%%%%%%%%%%%%%%%%%
%%%%%%%%%%%%%%%%%%%%%%%%%%%%%%%%%%%%%%%%%%%%%%%%%%%%

\section{Organização do Trabalho}
Para facilitar a leitura e a compreensão da pesquisa, este documento está estruturado da seguinte forma:

O Capítulo \ref{cap:conceitos} apresenta os aspectos conceituais que fundamentam o projeto. Nele, são abordados o funcionamento do protocolo LoRaWAN e o paradigma da criptografia leve (\textit{Lightweight Cryptography}), detalhando o processo de padronização do NIST que elegeu a família Ascon, além de uma comparação com outras famílias de algoritmos. O capítulo também explora a base teórica de \textit{Machine Learning} voltada para a detecção de anomalias.

O Capítulo \ref{cap:metodo} descreve o método do trabalho. Primeiramente, apresenta-se a 
abordagem de elicitação de requisitos adotada, incluindo análise de 
literatura, identificação de restrições de hardware e validação de 
viabilidade técnica. Em seguida, são detalhados os datasets analisados e 
selecionados para o modelo de Machine Learning, os procedimentos de 
treinamento e avaliação, bem como a justificativa para a adoção da 
biblioteca ascon-suite na implementação do padrão Ascon. O capítulo abrange 
ainda as etapas de configuração do ambiente, implementação e métodos de teste.

O Capítulo \ref{cap:especificacao} trata da especificação de requisitos do sistema, definindo as restrições computacionais, de memória e de consumo energético que o protótipo deve respeitar para operar na borda da rede.

O Capítulo \ref{cap:desenvolvimento} é dedicado ao projeto e desenvolvimento. Nesta etapa, são descritas as tecnologias e ferramentas práticas utilizadas, a arquitetura montada, os testes realizados e a discussão dos resultados obtidos.

Por fim, o Capítulo \ref{cap:consideracoes} apresenta as considerações finais, sumarizando as conclusões extraídas da pesquisa e apontando direcionamentos para futuras implementações e melhorias que podem surgir a partir deste trabalho.